\begin{lemma}
\label{editorial-source-lemma-primes-principal}
Let $R$ be a ring.
\begin{enumerate}
\item An ideal $I \subset R$ maximal with respect to not being
principal is prime.
\item If every prime ideal of $R$ is principal, then
every ideal of $R$ is principal.
\end{enumerate}
\end{lemma}

\begin{proof}
The first part follows from
Example \ref{algebra-example-oka-family-principal} and
Proposition \ref{algebra-proposition-oka}.
For the second, suppose that there exists an ideal $I \subset R$
which is not principal. In the nonempty set of nonprincipal ideals,
the empty chain has the chosen $I$ as an upper bound.
The union of a nonempty totally ordered chain
$\left\{I_\alpha\right\}$ of ideals that are not principal is not principal;
indeed, if $I = \bigcup I_\alpha$ were generated by
$a$, then $a$ would belong to some $I_\alpha $ and $a$ would generate it.
By Zorn's lemma, there is an ideal maximal with respect to not being
principal. This ideal is necessarily prime by the first part.
\end{proof}